% Railway single-track sequencing -- a big-M MILP.
% Standalone document: compile with  pdflatex problem.tex
\documentclass[11pt]{article}
\usepackage[a4paper,margin=2.4cm]{geometry}
\usepackage{amsmath,amssymb}
\usepackage{booktabs}
\usepackage{array}

\title{Railway Single-Track Sequencing\\[2pt]\large A big-M MILP}
\author{lp2graph example catalogue}
\date{}

\begin{document}
\maketitle

\section*{Problem}

A set of trains must each traverse one shared \emph{single-track segment}
(a tunnel, bridge, or passing-loop bottleneck). At most one train may
occupy the segment at a time, and a minimum \emph{headway} $h$ must elapse
between one train clearing the segment and the next entering it. Each train
$i$ has an earliest entry (release) time $r_i$, a running time $p_i$, and a
priority weight $w_i$. The dispatcher chooses each entry time $t_i$ (and
hence the order) to minimise the total weighted clearance time, where the
clearance time is $C_i = t_i + p_i$.

The mutual exclusion of any two trains $i,j$ -- \emph{either $i$ before $j$
or $j$ before $i$} -- is a disjunction, linearised with a binary ordering
variable and the \textbf{big-M} technique.

\section*{Sets, parameters, variables}

\begin{align*}
\mathcal{N} &= \{A,B,C,D\} && \text{trains}\\
\mathcal{P} &= \{(i,j)\in\mathcal{N}\times\mathcal{N} : i \prec j\} && \text{ordered index pairs}\\[4pt]
r_i,\ p_i,\ w_i && & \text{release, running time, weight of train }i\\
h && & \text{minimum headway}\\
M &= \max_{i} r_i + \textstyle\sum_{i\in\mathcal{N}} p_i + (|\mathcal{N}|-1)\,h && \text{big-M constant}\\[4pt]
t_i &\ge 0 && \text{entry time of train }i\\
C_i &\ge 0 && \text{clearance (completion) time of train }i\\
y_{ij} &\in\{0,1\} && y_{ij}=1 \iff \text{train }i\text{ enters before train }j,\ (i,j)\in\mathcal{P}
\end{align*}

\section*{Model}

\begin{align}
\min_{t,C,y}\quad & \sum_{i\in\mathcal{N}} w_i\,C_i \\
\text{s.t.}\quad
& t_i \ge r_i && \forall i\in\mathcal{N} \\
& C_i = t_i + p_i && \forall i\in\mathcal{N} \\
& t_j \ge t_i + p_i + h - M\,(1 - y_{ij}) && \forall (i,j)\in\mathcal{P} \\
& t_i \ge t_j + p_j + h - M\,y_{ij} && \forall (i,j)\in\mathcal{P} \\
& t_i \ge 0,\quad C_i \ge 0,\quad y_{ij}\in\{0,1\}.
\end{align}

\paragraph{Why big-M works.} For a pair $(i,j)$ the binary $y_{ij}$ selects
which of the two mutually exclusive sequencing constraints is active. If
$y_{ij}=1$, constraint (4) becomes the tight $t_j \ge t_i + p_i + h$ while
constraint (5) reads $t_i \ge t_j + p_j + h - M$, which for large $M$ is
vacuous. If $y_{ij}=0$ the roles swap. $M$ must be large enough never to cut
a feasible schedule, yet as small as possible so the LP relaxation stays
tight.

\section*{Data instance}

\begin{center}
\begin{tabular}{c r r r}
\toprule
train & $r_i$ & $p_i$ & $w_i$ \\
\midrule
A & 0 & 5 & 1 \\
B & 2 & 3 & 2 \\
C & 1 & 4 & 1 \\
D & 4 & 2 & 3 \\
\bottomrule
\end{tabular}
\end{center}

\noindent Headway $h = 1$, and
$M = \max_i r_i + \sum_i p_i + (|\mathcal{N}|-1)h = 4 + 14 + 3 = 21$.

\section*{Optimal solution}

The optimal weighted clearance is $\sum_i w_i C_i = 66$, attained by the
order $B \to D \to C \to A$:

\begin{center}
\begin{tabular}{c r r r}
\toprule
train & $t_i$ & $C_i$ & $w_i C_i$ \\
\midrule
B & 2  & 5  & 10 \\
D & 6  & 8  & 24 \\
C & 9  & 13 & 13 \\
A & 14 & 19 & 19 \\
\midrule
\multicolumn{3}{r}{total} & \textbf{66} \\
\bottomrule
\end{tabular}
\end{center}

\end{document}
